% Archivo generado automáticamente
\makeglossaries

\newglossaryentry{gnu}{
    name={gnu},
    description={},
    sort={gnu}
}

\newglossaryentry{a}{
    name={a},
    description={},
    sort={a}
}

\newglossaryentry{acetileno}{
    name={acetileno},
    description={Compuesto org\'anico de dos carbonos unidos mediante una triple ligadura},
    sort={acetileno}
}

\newglossaryentry{acido}{
    name={\'acido},
    description={(1) Sustancia donadora de iones H$^+$ (H$_3$O$^+$) (2) Sustancia que acepta electrones. (3) Al encontrarse en soluci\'on acuosa y estar contacto con metales produce sales},
    sort={acido}
}

\newglossaryentry{acidocarboxilico}{
    name={\'acido carboxilico},
    description={Sustancia org\'anica capaz de donar iones H$^+$ debido a que posee un grupo carbonilo},
    sort={acidocarboxilico}
}

\newglossaryentry{acidonucleico}{
    name={\'acido nucl\'eico},
    description={"\'Acidos org\'anicos que se encuentran en el n\'ucleo de las c\'elulas, son escenciales para la vida. Entre \'estos se encuentran el ADN y el ARN},
    sort={acidonucleico}
}

\newglossaryentry{acidograso}{
    name={{\'a}cidos grasos saturados},
    description={\'Acido carbox\'{\i}lico de cadena larga que usualmente se encuentra
en las grasas vegetales y animales},
    sort={acidograso}
}

\newglossaryentry{acidoorganico}{
    name={\'acido org\'anico},
    description={Son los \'acido carbox\'{\i}licos},
    sort={acidoorganico}
}

\newglossaryentry{adp}{
    name={Adenosin Difosfato},
    description={Sustancia compuesta de dos fosfatos que con\-tie\-ne enlaces de alta energ\'{\i}a},
    sort={adp}
}

\newglossaryentry{afinidadelectronica}{
    name={afinidad electr\'onica},
    description={Energ\'{\i}a desprendida o absorbida por un \'atomos al adicionarle un electr\'on, y es una medida de la atracci\'on que ejerce un \'atomo hacia un electr\'on},
    sort={afinidadelectronica}
}

\newglossaryentry{alcanos}{
    name={alcanos},
    description={Compuestos org\'anicos saturados con enlaces sencillos entre sus carbonos},
    sort={alcanos}
}

\newglossaryentry{alcoholes}{
    name={alcoholes},
    description={Compuestos org\'anicos con un grupo hidroxilo -OH},
    sort={alcoholes}
}

\newglossaryentry{alcoxido}{
    name={alc\'oxido},
    description={Compuesto i\'onico producto de la reacci\'on del sodio con un alcohol},
    sort={alcoxido}
}

\newglossaryentry{aldehido}{
    name={aldehido},
    description={Compuesto org\'anico que contiene un grupo carbonilo y al menos un hidr\'ogeno se encuentra enlazado a \'este},
    sort={aldehido}
}

\newglossaryentry{alifatico}{
    name={alif\'atico},
    description={Compuesto org\'anico de cadena abierta de \'atomos de carbono, puede contener ligaduras sencillas, dobles o triples},
    sort={alifatico}
}

\newglossaryentry{alqueno}{
    name={alquenos},
    description={Compuesto org\'anico con por lo menos un enlace doble entre sus carbonos},
    sort={alqueno}
}

\newglossaryentry{alquillitio}{
    name={alquil-litio},
    description={Compuesto org\'anomet\'alico producto de la reacc\'on de un haluro de alquilo con litio},
    sort={alquil}
}

\newglossaryentry{alquilo}{
    name={alquilo},
    description={Compuesto org\'anico con un \'atomo menos de hidr\'ogeno que en el alcano correspondiente},
    sort={alquilo}
}

\newglossaryentry{alquino}{
    name={alquinos},
    description={Compuesto org\'anico con  por lo menos un enlace triple entre sus carbonos},
    sort={alquino}
}

\newglossaryentry{amina}{
    name={aminas},
    description={Compuesto org\'anico que contiene el grupo NH$_x$ donde $x$ vale $0,1$ \'o $2$},
    sort={amina}
}

\newglossaryentry{anfotericos}{
    name={anfot\'ericos},
    description={Sustancias capaces de reaccionar tanto como \'acidos como bases},
    sort={anfoterico}
}

\newglossaryentry{antropica}{
    name={antr\'opica},
    description={},
    sort={antropica}
}

\newglossaryentry{arilo}{
    name={arilo},
    description={Es cualquier grupo funcional o sustituyente derivado de un anillo arom\'atico},
    sort={arilo}
}

\newglossaryentry{atp}{
    name={Adenosin Trifosfato},
    description={Sustancia compuesta de tres fosfatos que contiene enlaces de alta energ\'{\i}a},
    sort={atp}
}

\newglossaryentry{boyle}{
    name={Boyle},
    description={Robert Boyle estudio el comportamiento de los gases en funci\'on de la temperatura y presi\'on de manera cuantitativa},
    sort={boyle}
}

\newglossaryentry{BTU}{
    name={BTU},
    description={British Termal Unit. Calor necesario para incrementar 1$^\circ$F a una libra de agua},
    sort={btu}
}

\newglossaryentry{buffer}{
    name={buffer},
    description={Soluci\'on que amortigua los cambios en el $pH$ ya que se encuentra formada por \'acido y su base conjugada},
    sort={buffer}
}

\newglossaryentry{calor}{
    name={calor},
    description={Cantidad que fluye a trav\'es de la frontera de un sistema durante un cambio
de estado, en virtud de una diferencia de temperatura entre el sistema y su medio
exterior, y que fluye de un punto de mayor a uno de menor temperatura},
    sort={calor}
}

\newglossaryentry{caloria}{
    name={caloria},
    description={Cantidad de energ\'{\i}a necesaria para elevar en 1$^\circ$C a un gramo de agua},
    sort={caloria}
}

\newglossaryentry{calorimetro}{
    name={calor\'{\i}metro},
    description={Equipo empleado para medir la cantidad de calor aborbida o desprendiada por
un proceso f\'{\i}sico o qu\'{\i}mico},
    sort={calorimetro}
}

\newglossaryentry{capacidadcalorif}{
    name={capacidad calorifica},
    description={Cantidad de calor necesaria para producir un incremento de un grado en la
temperatura},
    sort={capacidadcalorif}
}

\newglossaryentry{carbonilo}{
    name={carbonilo},
    description={Grupo org\'anico que contiene un carbono y unido a \'este mediante una
 doble ligadura un ox\'{\i}geno, se encuentra presente en aldeh\'{\i}dos y cetonas},
    sort={carbonilo}
}

\newglossaryentry{catalizador}{
    name={catalizador},
    description={Es aquella sustancia que modifica la velocidad de reacci\'on y no sufre cambio
              significativo},
    sort={catalizador}
}

\newglossaryentry{cetona}{
    name={cetonas},
    description={Compuesto org\'anico que contiene un grupo carbonilo y unido a este dos grupos alquilo o arilo },
    sort={cetona}
}

\newglossaryentry{charles}{
    name={Charles},
    description={Estudio el comportamiento de los gases en funci\'on de la temperatura y el volumen de manera cuantitativa},
    sort={charles}
}

\newglossaryentry{constantedeacidez}{
    name={constante de acidez},
    description={Indica que tan disociado se encuentra un
\'acido, as\'{\i} como tambi\'en que tan fuerte es},
    sort={constantedeacidez}
}

\newglossaryentry{complejoactivado}{
    name={complejo activado},
    description={Especie intermedia entre reactivos y productos en una re\-ac\-ci\'on},
    sort={complejoactivado}
}

\newglossaryentry{craqueo}{
    name={craqueo},
    description={Deshidrogenaci\'on de alcanos por medio de altas temperaturas},
    sort={craqueo}
}

\newglossaryentry{deshidrogenacion}{
    name={deshidrogenaci\'on},
    description={Remoci\'on de hidr\'ogenos de una mol\'ecula org\'anica para formar dobles o triples enlaces entre \'atomos de carbono},
    sort={deshidrogenacion}
}

\newglossaryentry{deshidrohalogenacion}{
    name={deshidrohalogenaci\'on},
    description={Remoci\'on de hal\'ogenos de una mol\'ecula org\'anica para formar dobles o triples enlaces entre los \'atomos de carbono},
    sort={deshidrohalogenacion}
}

\newglossaryentry{disacaridos}{
    name={disac\'aridos},
    description={Compuestos por dos unidades de mon\'omeros},
    sort={disacaridos}
}

\newglossaryentry{ecuaciondeestado}{
    name={ecuaciones de estado},
    description={Representan las relaciones  $P$-$V$-$T$ de los
gases y describen completamente el estado de cualquier sistema gaseoso},
    sort={ecuaciondeestado}
}

\newglossaryentry{electronegatividad}{
    name={electronegatividad},
    description={Fuerza de atracci\'on que un \'atomo de un elemento presenta hacia los
               electrones en una mol\'ecula},
    sort={electronegatividad}
}

\newglossaryentry{endergonicas}{
    name={enderg\'onica},
    description={Reacciones bioqu\'{\i}micas que para ocurrir requieren de energ\'{\i}a para
efectuarse},
    sort={endergonicas}
}

\newglossaryentry{endotermico}{
    name={endot\'ermico},
    description={Proceso que absorbe energ\'{\i}a al efectuarse},
    sort={endotermico}
}

\newglossaryentry{energiadeactivacion}{
    name={energ\'{\i}a de activaci\'on},
    description={Energ\'{\i}a necesaria para iniciar una reacci\'on qu\'{\i}mica},
    sort={energiadeactivacion}
}

\newglossaryentry{entalpia}{
    name={entalpia},
    description={Mide los cambios t\'ermicos a presi\'on constante en un sistema, los cambios
pueden ser qu\'{\i}micos (entalp\'{i}a de reacci\'on, combusti\'on, formaci\'on) o
f\'{\i}sicos entalp\'{\i}a de disoluci\'on)},
    sort={entalpia}
}

\newglossaryentry{entropia}{
    name={entropia},
    description={Grado de desorden en un sistema. Se mide en unidades en\-tr\'o\-pi\-cas ue
cal/grado},
    sort={entropia}
}

\newglossaryentry{exergonicas}{
    name={exerg\'onica},
    description={Reacciones bioqu\'{\i}micas que almacenan energ\'{\i}a en enlaces de alta
energ\'{\i}a al efectuarse},
    sort={exergonicas}
}

\newglossaryentry{exotermico}{
    name={exot\'ermica},
    description={Proceso que desprende energ\'{\i}a al efectuarse},
    sort={exotermico}
}


\newglossaryentry{eter}{
    name={{\'e}ter},
    description={Es un compuesto org\'anico que contiene dos grupos alquilo o arilo (grupo \'eter) con un \'atomo de ox\'{\i}geno},
    sort={eter}
}

\newglossaryentry{egne}{
    name={egne},
    description={},
    sort={egne}
}

\newglossaryentry{ester}{
    name={éster},
    description={Es un compuesto  de un radical org\'anico unido al residuo de cualquier \'acido oxigenado (ox\'acido ), org\'anico o inorg\'anico. },
    sort={ester}
}

\newglossaryentry{factordecompresibilidad}{
    name={factor de compresibilidad},
    description={Es la raz\'on del volumen molar observado con respecto al volumen molar ideal en funci\'on de la presi\'on a temperatura constante $Z=\frac{p\overline V}{RT}$},
    sort={factordecompresibilidad}
}

\newglossaryentry{fase}{
    name={fase},
    description={Parte de un sistema que es fisicamente diferenciable y se encuentra separado de las otras partes del sistema por un l\'imite. ej. un sistema de dos fases es agua con aceite},
    sort={fase}
}

\newglossaryentry{fenilo}{
    name={fenilo},
    description={Compuesto derivado del benceno},
    sort={fenilo}
}

\newglossaryentry{frontera}{
    name={frontera},
    description={Termodin\'amicamente hablando es l\'{\i}mite real o imaginario que define
a un sistema},
    sort={frontera}
}

\newglossaryentry{fem}{
    name={fem},
    description={Fuerza electromotriz},
    sort={fem}
}

\newglossaryentry{gasesideales}{
    name={gases ideales},
    description={Son aquellos en el que su comportamiento no se ve influenciado por las
fuerzas de cohesi\'on y vol\'umenes de las mol\'eculas que los componen},
    sort={gasesideales}
}

\newglossaryentry{gasesnobles}{
    name={gases nobles},
    description={Son aquellos que  tienen completa la capa de valencia con 8 electrones, tienen propiedades similares como la baja reactividad  y son  monoat\'omicos en condiciones normales},
    sort={gasesnobles}
}

\newglossaryentry{gasesreales}{
    name={gases reales},
    description={Poseen un comportamiento similar al de los ideales cuando se encuentran a temperaturas altas y presiones bajas},
    sort={gasesreales}
}

\newglossaryentry{genotoxicas}{
    name={genot\'oxicas},
    description={Son gentes f\'{\i}sicos (temperatura, luz ultravioleta, radiaciones ionizantes, radiaciones electromagn\'eticas, etc.) o productos qu\'{\i}micos capaces de alterar la informaci\'on gen\'etica celular},
    sort={genotoxicas}
}

\newglossaryentry{halogenacion}{
    name={halogenaci\'on},
    description={Proceso de adicionar hal\'ogenos a una mol\'ecula},
    sort={halogenacion}
}

\newglossaryentry{helicoidal}{
    name={helicoidal},
    description={Que tiene forma de h\'elice},
    sort={helicoidal}
}

\newglossaryentry{Henry}{
    name={Henry},
    description={La Ley de Henry es una ley que describe la relaci\'on entre la presi\'on y la solubilidad de los gases en soluciones},
    sort={henry}
}

\newglossaryentry{heteropolimero}{
    name={heteropolimero},
    description={Pol\'{\i}mero compuesto de dos o m\'as mo\-n\'o\-meros diferentes},
    sort={heteropolimero}
}

\newglossaryentry{hidrocarburos}{
    name={hidrocarburos},
    description={Compuestos constituidos por \'atomos de carbono e hidr\'ogeno unidos entre s\'{\i} mediante enlaces covalentes},
    sort={hidrocarburos}
}

\newglossaryentry{hidrolisis}{
    name={hidr\'olisis},
    description={Reacci\'on de una sustancia i\'onica con el agua},
    sort={hidrolisis}
}

\newglossaryentry{hibridacion}{
    name={hibridaci\'on},
    description={Combinaci\'on de los orbitales $s$ y $p$ del \'atomo de carbono},
    sort={hibridacion}
}

\newglossaryentry{homopolimero}{
    name={homopolimero},
    description={Pol\'{\i}mero compuesto de un solo tipo de mo\-n\'o\-mero},
    sort={homopolimero}
}

\newglossaryentry{ionizacion}{
    name={ionizaci\'on},
    description={La cantidad de energ\'{\i}a necesaria para quitar un electr\'on a un \'atomo},
    sort={ionizacion}
}

\newglossaryentry{ionhidronio}{
    name={i\'on hidronio},
    description={Combinaci\'on de un prot\'on con una mol\'ecula polar dando un i\'on
hidratado H$_3$O$^+$},
    sort={ionhidronio}
}

\newglossaryentry{isomeros}{
    name={is\'omeros},
    description={Compuestos qu\'{\i}micos diferentes con la misma f\'ormula condensada},
    sort={isomeros}
}

\newglossaryentry{meteorizacion}{
    name={meteorizaci\'on},
    description={El proceso de desintegraci\'on o descomposici\'on de las rocas y minerales en la superficie terrestre o cerca de ella},
    sort={meteorizacion}
}

\newglossaryentry{mie}{
    name={Mie},
    description={Esparcimiento que se produce cuando la luz choca con part\'{\i}culas o mol\'eculas grandes. Las part\'{\i}culas absorben una parte de la luz y reflejan el resto, como peque\~nos espejos},
    sort={mie}
}

\newglossaryentry{monosacaridos}{
    name={monosac\'aridos},
    description={Hidrocarburos compuestos por un mon\'omero},
    sort={monosacaridos}
}

\newglossaryentry{orbitales}{
    name={orbitales},
    description={Los principales niveles de energ\'{\i}a donde se localizan los electrones en
               un \'atomo},
    sort={orbitales}
}

\newglossaryentry{oligosacaridos}{
    name={oligosac\'aridos},
    description={Compuestos de 3 a 10 unidades de mon\'omeros},
    sort={oligosacaridos}
}

\newglossaryentry{pargalvanico}{
    name={par galv\'anico},
    description={Contacto entre dos metales de diferente electronegatividad en
condiciones favorables es posible la transferencia de electrones},
    sort={pargalvanico}
}

\newglossaryentry{ph}{
   type=symbols,
    name={$pH$},
    description={Logaritmo negativo de la concentraci\'on de iones H$_3$O$^+$, nos indica
que tan \'acida ($pH < 7$) o alcalina ($p H >7$) es una sustancia},
    sort={ph}
}

\newglossaryentry{polaridad}{
    name={polaridad},
    description={Es el grado al cual un par electr\'onico es desigualmente compartido en una
mol\'ecula},
    sort={polaridad}
}

\newglossaryentry{pirolisis}{
    name={pir\'olisis},
    description={Deshidrogenaci\'on de alcanos a altas temperaturas de entre 500 a 700
$^\circ$C},
    sort={pirolisis}
}

\newglossaryentry{polipeptidos}{
    name={polip\'eptidos},
    description={Cadenas de entre 40 a 50 amino\'aciodos unidos por enlces
pep\-t\'{\i}\-di\-cos},
    sort={polipeptidos}
}

\newglossaryentry{poliprotico}{
    name={polipr\'otico},
    description={Sustancia que contribuye con m\'as de un prot\'on al agua},
    sort={poliprotico}
}

\newglossaryentry{polisacaridos}{
    name={polisac\'aridos},
    description={Carbohidratos compuestos de m\'as de 10 unidades de mono\-sa\-c\'a\-ri\-dos},
    sort={polisacaridos}
}

\newglossaryentry{sistema}{
    name={sistema},
    description={Una porci\'on del universo que nos interesa estudiar y se encuentra separada
por fronteras},
    sort={sistema}
}

\newglossaryentry{solvente}{
    name={solvente},
    description={Sustancia l\'{\i}quida en mayor proporci\'on en una soluci\'on},
    sort={solvente}
}

\newglossaryentry{soluto}{
    name={soluto},
    description={Sustancia en que se encuentra en menor proporci\'on en una soluci\'on},
    sort={soluto}
}

\newglossaryentry{temperaturaabsoluta}{
    name={temperatura absoluta},
    description={Temperatura a la cual las mol\'eculas de un gas se encuentran en forma
cristalina y sin ning\'un tipo de movimiento},
    sort={temperaturaabsoluta}
}

\newglossaryentry{termodinamica}{
    name={termodin\'amica},
    description={Estudia las diferentes clases de energ\'{\i}a y sus diferentes ma\-ni\-festaciones},
    sort={termodinamica}
}

\newglossaryentry{termoestables}{
    name={termoestables},
    description={Pol\'{\i}meros que no llegan a fundirse al calentarlos},
    sort={termoestables}
}

\newglossaryentry{termoplasticos}{
    name={termopl\'asticos},
    description={Pol\'{\i}meros que se ablandan con el aumento gradual de la temperatura y
llegan a fundirse},
    sort={termoplasticos}
}

\newglossaryentry{tioeter}{
    name={tio\'eter},
    description={Un compuesto  similar a un \'eter, conteniendo un \'atomo de azufre en vez de un \'atomo de ox\'{\i}geno},
    sort={tioeter}
}

\newglossaryentry{tiol}{
    name={tiol},
    description={Un compuesto que contiene el grupo funcional formado por un \'atomo de azufre y un \'atomo de hidr\'ogeno (-SH)},
    sort={tiol}
}

\newglossaryentry{trabajo}{
    name={trabajo},
    description={Termodin\'amicamente hablando trabajo es la cualquier cantidad de
energ\'{\i}a qur fluye a trav\'es de la frontera de un sistema durante un cambio de
estado y que puede ser empleada para elevar un cuerpo en el medio exterior},
    sort={trabajo}
}

\newglossaryentry{UIQPA}{
    name={UIQPA },
    description={Uni\'on Internacional de Qu\'{\i}mica Pura y Aplicada (International Union of Pure and Applied Chemistry},
    sort={UIQPA}
}

\newglossaryentry{valencia}{
    name={valencia},
    description={Los electrones en la capa externa de un \'atomo que son los responsables de la
mayor parte de la actividad electr\'onica},
    sort={valencia}
}

\newglossaryentry{variabledeestado}{
    name={variable de estado},
    description={Es aquella que tiene un valor definido cuando se especifica el estado de un
sistema},
    sort={variabledeestado}
}

\newglossaryentry{vulcanizacion}{
    name={vulcanizaci\'on},
    description={Adici\'on de azufre al hule natural caliente para modificar as\'{\i} sus
propiedades},
    sort={vulcanizacion}
}

\newglossaryentry{nuname}{
    name={nu},
    description={Nombre de la letra griega $\nu$},
    sort={nuname}
}

\newglossaryentry{exclam}{
    name={!},
    description={},
    sort={exclam}
}

\newglossaryentry{interr}{
    name={?},
    description={},
    sort={interr}
}

\newglossaryentry{ASCII}{
    name={American Standard Code for Information Interchange},
    description={El sistema m\'as popular de codificaci\'on},
    sort={ASCII}
}

\newglossaryentry{bit}{
    name={bit},
    description={},
    sort={bit}
}

\newglossaryentry{bibtex8}{
    name={bib\TeX8},
    description={},
    sort={bibtex8}
}

\newglossaryentry{pi}{
   type=symbols,
    name={$\pi$},
    description={Constante matem\'atica aproximadamente igual a la 
  versi\'on m\'as reciente de \TeX},
    sort={pi}
}

\newglossaryentry{alfa}{
    name={$\alpha$},
    description={},
    sort={alfa}
}

\newglossaryentry{deltae}{
   type=symbols,
    name={$\Delta {\mathrm E}$},
    description={cambio en la energ\'{\i}a del sistema},
    sort={deltae}
}

\newglossaryentry{deltav}{
   type=symbols,
    name={$\Delta {\mathrm V}$},
    description={cambio en el volumen del sistema},
    sort={deltav}
}

\newglossaryentry{deltahc}{
   type=symbols,
    name={$\Delta H_{\mathrm c}$},
    description={Entalp\'{\i}a de combusti\'on},
    sort={deltahc}
}

\newglossaryentry{deltahd}{
   type=symbols,
    name={$\Delta H_{\mathrm d}$},
    description={Entalp\'{\i}a de disoluci\'on},
    sort={deltahd}
}

\newglossaryentry{deltahe}{
  type=symbols,
    name={$\Delta {\mathrm H}_{25^\circ C}$},
    description={Entalp\'{\i}a de enlace a $25^\circ\, C$},
    sort={deltahe}
}

\newglossaryentry{deltahen}{
   type=symbols,
    name={$\Delta H_{\mathrm e}$},
    description={Entalp\'{\i}a de enlace},
    sort={deltahen}
}

\newglossaryentry{nu}{
   type=symbols,
    name={$\nu$},
    description={Frecuencia de la luz incidente},
    sort={nu}
}

\newglossaryentry{deltahev}{
   type=symbols,
    name={$\Delta H_{\mathrm b}$},
    description={Entalp\'{\i}a de vaporizaci\'on},
    sort={deltahev}
}

\newglossaryentry{deltahfr}{
   type=symbols,
    name={$\Delta H_{\mathrm f}$},
    description={Entalp\'{\i}a de formaci\'on},
    sort={deltahfr}
}

\newglossaryentry{deltahf}{
   type=symbols,
    name={$\Delta H_{\mathrm s}$},
    description={Entalp\'{\i}a de fusi\'on},
    sort={deltahf}
}

\newglossaryentry{deltag}{
   type=symbols,
    name={$\Delta {\mathrm G}$},
    description={Energ\'{\i}a libre de Gibbs},
    sort={deltag}
}

\newglossaryentry{deltahr}{
   type=symbols,
    name={$\Delta H_{\mathrm r}$},
    description={Entalp\'{\i}a de reacci\'on},
    sort={deltahr}
}

\newglossaryentry{deltas}{
   type=symbols,
    name={$\Delta {\mathrm S}$},
    description={Cambio de entrop\'{\i}a},
    sort={deltas}
}

\newglossaryentry{densidad}{
   type=symbols,
    name={$\rho$},
    description={Densidad del gas o l\'{\i}quido en masa/volumen [\kilogrampercubicmetrenp]},
    sort={densidad}
}

\newglossaryentry{ecelda}{
   type=symbols,
    name={$E_\textrm{celda}$},
    description={Potencial de celda, [volt]},
    sort={ecelda}
}

\newglossaryentry{ereducc}{
   type=symbols,
    name={$E_\textrm{celda}$},
    description={Potencial de oxidaci\'on, [volt]},
    sort={ereducc}
}

\newglossaryentry{eoxida}{
   type=symbols,
    name={$E_\textrm{celda}$},
    description={Constante de Planck [\joule \reciprocal\second]},
    sort={eoxida}
}

\newglossaryentry{hplank}{
   type=symbols,
    name={$h$},
    description={Constante de Planck [\joule \reciprocal\second]},
    sort={hplank}
}

\newglossaryentry{zeta}{
    name={z},
    description={},
    sort={zeta}
}

\newglossaryentry{awaals}{
   type=symbols,
    name={$a$},
    description={Constante de la ecuaci\'on de Van der Waals},
    sort={awaals}
}

\newglossaryentry{bwaals}{
   type=symbols,
    name={$b$},
    description={Constante de la ecuaci\'on de Van der Waals},
    sort={bwaals}
}

\newglossaryentry{cp}{
   type=symbols,
    name={$C_p$},
    description={Capacidad calor\'{\i}fica a presi\'on constante [cal/(g K)]},
    sort={cp}
}

\newglossaryentry{H}{
   type=symbols,
    name={$H$},
    description={Entalp\'{\i}a},
    sort={H}
}

\newglossaryentry{extincion}{
   type=symbols,
    name={$\sigma$},
    description={Coeficiente de extinci\'on},
    sort={extincion}
}

\newglossaryentry{ggravedad}{
   type=symbols,
    name={$g$},
    description={Aceleraci\'on debida a la gravedad, 9.81 m/s$^2$},
    sort={ggravedad}
}

\newglossaryentry{haltura}{
    name={$h$},
    description={Distancia en altura, [m]},
    sort={haltura}
}

\newglossaryentry{gamma}{
   type=symbols,
    name={$\Gamma$},
    description={Gradiente adiab\'atico seco},
    sort={gamma}
}

\newglossaryentry{londa}{
   type=symbols,
    name={$\lambda$},
    description={Longitud de onda de la luz},
    sort={londa}
}

\newglossaryentry{masa}{
   type=symbols,
    name={$m$},
    description={Masa de un cuerpo, [kg]},
    sort={masa}
}

\newglossaryentry{moles}{
   type=symbols,
    name={$n$},
    description={N\'umero de moles [gmol] $=$ peso[g]/peso molecular[g/gmol]},
    sort={moles}
}

\newglossaryentry{niveles}{
   type=symbols,
    name={n},
    description={N\'umero de niveles de energ\'{\i}a en un \'atomo},
    sort={niveles}
}

\newglossaryentry{subnivels}{
   type=symbols,
    name={$s$},
    description={Subnivel de energ\'{\i}a en un \'atomo consta de un orbital},
    sort={subnivels}
}

\newglossaryentry{subnivelp}{
   type=symbols,
    name={$p$},
    description={Subnivel de energ\'{\i}a en un \'atomo consta de tres orbitales},
    sort={subnivelp}
}

\newglossaryentry{subniveld}{
   type=symbols,
    name={$d$},
    description={Subnivel de energ\'{\i}a en un \'atomo consta de cinco orbitales},
    sort={subniveld}
}

\newglossaryentry{subnivelf}{
   type=symbols,
    name={$f$},
    description={Subnivel de energ\'{\i}a en un \'atomo consta de siete orbitales},
    sort={subnivelf}
}

\newglossaryentry{tresid}{
   type=symbols,
    name={$\tau$},
    description={Tiempo de residencia o vida \'util},
    sort={tresid}
}

\newglossaryentry{presion}{
   type=symbols,
    name={$P$},
    description={Presi\'on en un sistema},
    sort={presion}
}

\newglossaryentry{qcalor}{
   type=symbols,
    name={$Q$},
    description={Energ\'{\i}a en foma de calor},
    sort={qcalor}
}

\newglossaryentry{constR}{
   type=symbols,
    name={$R$},
    description={Constante de la ley general de los gases, 0.082057
$\frac{l-atm}{gmol\;K}$},
    sort={constR}
}

\newglossaryentry{temperatura}{
   type=symbols,
    name={$T$},
    description={Temperatura en un sistema, [K]},
    sort={temperatura}
}

\newglossaryentry{orbitalsp}{
   type=symbols,
    name={$sp$},
    description={Orbital h\'{\i}brido del carbono entre orbitales $s$ y $p$},
    sort={orbitalsp}
}

\newglossaryentry{volumen}{
   type=symbols,
    name={$V$},
    description={Volumen en un sistema},
    sort={volumen}
}

\newglossaryentry{wtrab}{
   type=symbols,
    name={$w$},
    description={Trabajo realizado por el sistema sobre el contorno},
    sort={wtrab}
}

\newglossaryentry{fcomp}{
   type=symbols,
    name={$z$},
    description={Factor de compresibilidad},
    sort={fcomp}
}

